\documentclass[12pt]{article}
\usepackage[margin=1in]{geometry}
\usepackage{hyperref}
\usepackage{parskip}   % space between paragraphs, no indent

\begin{document}

%% ── Sender block ──────────────────────────────────────────────────────────────
\noindent
Eden Ofek Nalian \\
Department of Mathematics \& Statistics \\
McMaster University \\
Hamilton, Ontario, Canada \\
\href{mailto:eden.nalian@gmail.com}{eden.nalian@gmail.com}

\bigskip
\noindent \today

\bigskip

%% ── Recipient block ───────────────────────────────────────────────────────────
\noindent
Editorial Office \\
\textit{PLOS Digital Health}

\bigskip

%% ── Salutation ────────────────────────────────────────────────────────────────
\noindent Dear Editors,

\bigskip

%% ── Body ──────────────────────────────────────────────────────────────────────
We are writing to submit our manuscript, \textbf{``Validation of automated
table extraction for historical infectious-disease surveillance data,''} for consideration as a
Research Article in \textit{PLOS Digital Health}.


The majority of historical public health data remain locked in scanned paper
documents, representing a major barrier to large-scale epidemiological research.
Our manuscript addresses this barrier directly by rigorously evaluating
Amazon Textract (a cloud-based machine-learning service for structured
document extraction) as an  alternative to manual data entry for historical infectious disease tables.
We apply the tool to the Canadian Disease Incidence Dataset (CANDID), an archive of weekly  disease counts spanning multiple Canadian provinces and territories from 1956 to 1958.
This work sits squarely within \textit{PLOS Digital Health}'s scope: it uses
a digital tool to advance the accessibility and usability of health data, and it provides a transparent, reproducible methodology evaluation with direct implications for public health research practice.

After a straightforward post-processing pipeline, automated extraction
achieved over 95\% cell-level agreement with manually entered data.
The extraction errors that remained were either small in magnitude and statistically negligible, or large and visually obvious, making them easy to detect and correct.
We further demonstrate that epidemiologically important statistics derived from the extracted data (including disease cross-correlations, epidemic peak times, peak case counts, and exponential growth rates) were nearly
identical to those computed from the manually entered dataset. This establishes
that automated table extraction can produce data of sufficient quality for
substantive epidemiological analysis.
We also identify conditions under which automation is most beneficial and
outline a clear post-processing workflow that other research groups can adopt
for their own historical archives.


The bottleneck of manual data entry has long constrained historical
epidemiology, forcing researchers to focus on small datasets or invest
enormous person-hours before any analysis can begin.
As archives of historical health records continue to be scanned worldwide,
tools and protocols like Amazon Textract   will become increasingly important for realising the full scientific value of these
collections. We believe this paper will be of broad interest to readers working in
digital health informatics, public health surveillance, mathematical
epidemiology, and data science.

We confirm that this manuscript has not been published elsewhere and is
not under consideration at any other journal.
All authors have read and approved the final manuscript.
There are no competing interests to declare, and the study received no
external funding.

We have not indicated a preferred Academic Editor, and we leave the
assignment entirely to the editorial team's discretion.
Suggested reviewers are listed in the submission form.

We look forward to hearing from you.

\bigskip

\noindent Sincerely,

\bigskip
\bigskip

\noindent Eden Ofek Nalian (corresponding author) \\
Steven C.\ Walker and David J.\,D.\ Earn  

\end{document}